\section{Guida all'esecuzione}
\label{sec:esecuzione}

Questa sezione fornisce le istruzioni operative per eseguire il progetto.
L'archivio consegnato include già tutto il necessario per l'esecuzione
immediata: il checkpoint della rete allenata
(\texttt{models/CNN\_64f-k3-3blk.pth}) e il sottoinsieme di test già
estratto (\texttt{plantvillage\_testset/}, Sezione~\ref{sec:moduli}).
Senza questi due elementi non sarebbe possibile eseguire un'inferenza
sensata: il primo perché il training della rete richiede ore su GPU e non è
lo scopo di questo progetto, il secondo perché senza di esso la fotocamera
virtuale pescherebbe anche immagini già viste in addestramento, mostrando
un'accuratezza artificiosamente vicina al 100\%.

\subsection{Prerequisiti}

\begin{itemize}
    \item Python 3.10 o superiore.
    \item I file già inclusi nell'archivio: checkpoint della rete e
    cartella \texttt{plantvillage\_testset/} (si veda sopra).
\end{itemize}

\subsection{Passo 1 -- Ambiente virtuale e dipendenze}

\begin{lstlisting}[language=bash]
# Crea un ambiente virtuale isolato
python3 -m venv venv

# Attiva l'ambiente
source venv/bin/activate          # macOS / Linux
# venv\Scripts\activate           # Windows

# Installa le dipendenze del progetto
pip install --upgrade pip
pip install -r requirements.txt
\end{lstlisting}

\subsection{Passo 2 -- Configurazione}

Verificare che il file \texttt{config.yaml} punti alla cartella del test
set già inclusa nell'archivio:

\begin{lstlisting}[language={}]
paths:
  dataset_root: "./plantvillage_testset"
  checkpoint: "./models/CNN_64f-k3-3blk.pth"
\end{lstlisting}

\begin{warnbox}
Se si utilizza la dashboard Streamlit, il file \texttt{config.yaml} viene
letto una sola volta all'avvio del processo e tenuto in cache per tutta la
sua durata: ogni modifica al file richiede l'arresto (\texttt{CTRL+C}) e il
riavvio del comando \texttt{streamlit run} -- un semplice aggiornamento
della pagina nel browser \emph{non} è sufficiente.
\end{warnbox}

\emph{Nota per la riproducibilità}: la cartella \texttt{plantvillage\_testset/}
è stata generata con lo script \texttt{scripts/extract\_test\_set.py}, che
riproduce lo split casuale (\texttt{seed=42}, 80/20) del progetto di
training originale a partire dal dataset completo. Rieseguirlo non è
necessario per usare il progetto così come consegnato; serve solo se si
vuole rigenerare il sottoinsieme di test da una copia locale del dataset
completo.

\subsection{Passo 3 -- Generazione dei modelli compressi (opzionale)}

Necessario solo per usare la modalità \texttt{optimized} o i grafici di
confronto nella dashboard:

\begin{lstlisting}[language=bash]
python models/compress.py
\end{lstlisting}

Al termine, \texttt{models/optimized/} conterrà le tre varianti compresse
(Sezione~\ref{sec:compressione}) e \\ \texttt{benchmarks/benchmark\_results.csv}
verrà aggiornato con le relative metriche.

\subsection{Passo 4 -- Avvio del sistema}

Le due modalità (Sezione~\ref{sec:moduli}) rispondono a esigenze diverse: in
un deployment reale, il Raspberry Pi nella serra girerebbe senza schermo,
scrivendo solo un log, cioè lo scenario riprodotto dall'esecuzione da riga di
comando; la dashboard, invece, rappresenta uno strumento di supervisione
che un operatore userebbe per ispezionare il sistema da remoto o durante
una demo, non qualcosa che girerebbe stabilmente sul dispositivo stesso.

\subsubsection{Dashboard interattiva (consigliata per la dimostrazione)}

\begin{lstlisting}[language=bash]
streamlit run dashboard/app_streamlit.py
\end{lstlisting}

Apre automaticamente il browser su \url{http://localhost:8501}. Dalla barra
laterale si scelgono modalità/variante del modello, un'eventuale classe
forzata per demo mirate, l'intervallo fra i cicli, e si avvia/ferma/resetta
la simulazione.

\subsubsection{Esecuzione da riga di comando}

\begin{lstlisting}[language=bash]
# Esecuzione base, 50 cicli
python orchestrator/main_loop.py --cycles 50 --interval 2

# Forzare una classe specifica (demo mirata)
python orchestrator/main_loop.py --class Tomato___Early_blight --cycles 20

# Selezionare esplicitamente modalità e variante del modello
python orchestrator/main_loop.py --mode optimized --variant onnx --cycles 100

# Esecuzione infinita (fino a CTRL+C), output minimo
python orchestrator/main_loop.py --cycles 0 --quiet
\end{lstlisting}

Il file di log generato in \texttt{logs/} può essere analizzato con
strumenti standard per l'analisi tabellare (es. \texttt{pandas}, o un
foglio di calcolo).
